\begin{table}[H]
\centering
\caption{Representative perturbation, counterfactual, and corruption-based evaluation studies in the dataset corpus.}
\label{tab:perturbation_methods}
\renewcommand{\arraystretch}{1.2}
\begin{tabular}{p{3.0cm} p{5.6cm} p{5.4cm}}
\hline
\textbf{Reference} & \textbf{Method} & \textbf{Application} \\
\hline
\cite{Zhang2021} 2021 & Input perturbation and attribution comparison & ECG diagnosis \\
\cite{MorenoSanchez2023} 2023 & Explainable feature ablation & Heart-failure survival \\
\cite{singla2023explaining} 2023 & Counterfactual explanation generation & Medical image analysis \\
\cite{lenis2020domain} 2020 & Counterfactual impact analysis & Medical image classifier interpretation \\
\cite{palatnik2021xai} 2021 & Bias-sensitive saliency analysis & COVID CT classifier auditing \\
\cite{Pertuz2023} 2023 & Saliency-mask comparison & Breast lesion detection \\
\cite{Repetto2025EvaluatingTR} 2025 & Corruption-based robustness testing & Medical image recognition \\
\cite{Patricio2024} 2024 & AOPC/MoRF perturbation evaluation of attribution maps & Medical image classification \\
\hline
\end{tabular}
\end{table}
